%==============================================================================
% sections/07_estimators_and_efficiency.tex
%==============================================================================
\section{Estimator Families and Efficiency Criteria}
\label{sec:estimators-efficiency}

\subsection{Estimator menu}

The expanded paper compares six estimator families.

\begin{description}
\item[Crude Monte Carlo.] Simulate \(X\) and average \(\1\{S_N>x\}\). This is
the baseline for sample-size and work-normalized comparisons.

\item[Independent summed CdMC.] Use \(\sum_i\Fb_i(T_i)\). This is exact only
under independence but remains a useful misspecification benchmark.

\item[Dependent kernel CdMC.] Use \(\sum_i\Pb(X_i>T_i\mid X_{-i})\). This is
exact under arbitrary dependence if conditional kernels are correct.

\item[Latent-shock CdMC.] Transform to independent shocks and apply independent
CdMC to signed shock contributions.

\item[Block CdMC.] Estimate or model block-tail probabilities and apply CdMC to
independent block sums.

\item[Spectral/radial CdMC.] Condition on angular direction and evaluate radial
survival. This is the preferred estimator when MRV with non-axis spectral mass
is empirically supported.
\end{description}

\begin{table}[p]
\centering
\begin{tabularx}{\linewidth}{p{0.18\linewidth}p{0.19\linewidth}p{0.22\linewidth}p{0.18\linewidth}X}
\toprule
Estimator & Conditioning object & Required model input & Exactness condition & Planned diagnostic \\
\midrule
Independent CdMC & Observed coordinate & Marginal survival functions & Observed contributions independent & Bias under dependence; BRE replication \\
Dependent kernel CdMC & Observed coordinate given others & Conditional survival kernels & Conditional kernels correctly specified & Kernel calibration and relative error \\
Latent-shock CdMC & Independent shock & Shock mixing matrix and shock tails & Shock representation correctly specified & Shock attribution and common-shock bias reduction \\
Block CdMC & Block sum & Block partition and block-tail model & Blocks independent or weakly coupled & Block stability and work-normalized error \\
Copula CdMC & Conditional copula coordinate & Margins plus conditional copula & Copula and margins correctly specified & Tail-dependence fit and backtest performance \\
Spectral CdMC & Angular direction & Spectral sample and radial tail & MRV radial-angular model valid & Spectral mass stability and radial tail fit \\
Hidden-cone mixture & Axis/hidden stratum & Hidden-cone scale and sampler & Hidden scale estimated reliably & Hidden term versus second-order scales \\
\bottomrule
\end{tabularx}

\caption[Dependent estimator summary]{Dependent estimator summary. The table
specifies the conditioning object, required model input, exactness condition,
and planned efficiency diagnostic for each estimator family.}
\label{tab:dependent-estimator-summary}
\end{table}

\subsection{Efficiency definitions}

We collect the efficiency notions used throughout the paper. Throughout,
\(Z(x)\) denotes a generic unbiased estimator of \(\mu(x)\), \(n\) the
replication count, and \(\mathrm{cost}(Z)\) the average wall-clock cost per
replicate (or equivalently, the count of tail-function evaluations per
replicate when wall-clock work is dominated by tail evaluations).

\begin{definition}[Bounded relative error, vanishing relative error, log-efficiency]
\label{def:efficiency-notions}
An unbiased estimator $Z(x)$ for $\mu(x)$:
\begin{itemize}
\item has \emph{bounded relative error} (BRE) if
  $\limsup_{x\to\infty}\Var Z(x)/\mu(x)^2<\infty$;
\item has \emph{vanishing relative error} (VRE) if
  $\Var Z(x)/\mu(x)^2\to 0$;
\item is \emph{logarithmically efficient} if
  $\liminf_{x\to\infty}\log\E Z(x)^2/\log\mu(x)^2\ge 1$.
\end{itemize}
\end{definition}

\begin{definition}[Work, work-normalized variance, efficiency frontier]
\label{def:work-normalized}
The \emph{work} of $Z(x)$ at threshold $x$ is the mean cost per replicate
$\mathrm{cost}(Z;x)\in(0,\infty)$. The \emph{work-normalized variance} is
\[
  W(Z;x)=\Var Z(x)\cdot\mathrm{cost}(Z;x),
\]
the \emph{work-normalized relative error} (WNRE) is
\[
  \mathrm{WNRE}(Z;x)=\sqrt{W(Z;x)}/\mu(x),
\]
and the \emph{efficiency ratio} of $Z_1$ relative to $Z_2$ is
$\mathrm{eff}(Z_1,Z_2;x)=W(Z_2;x)/W(Z_1;x)$. The \emph{efficiency frontier}
over a class of estimators $\mathcal{Z}$ at threshold $x$ is the lower convex
boundary of the set $\{(\mathrm{cost}(Z;x),\Var Z(x)):Z\in\mathcal{Z}\}$.
\end{definition}

In the empirical pipeline $\widehat{c}_{\mathrm{cpu}}$ denotes the sample
estimate of $\mathrm{cost}(Z;x)$ and $\widehat{\mathrm{WNRE}}(Z;x)$ replaces
the population quantities by their sample analogues.

\subsection{Efficiency results for the principal estimators}
\label{subsec:efficiency-propositions}

The independent and dependent CdMC efficiency results proved in
\cref{appx:independent-proofs,appx:dependent-proofs} are summarized here as
formal propositions for ease of reference.

\begin{proposition}[BRE of independent CdMC]
\label{prop:ind-cdmc-bre}
Under \cref{ass:ind-pte,ass:ind-sbj}, the independent summed CdMC estimator
$Z^{\mathrm{ind}}(x)$ is unbiased for $\mu(x)$ and satisfies
\(
  \limsup_{x\to\infty}\Var Z^{\mathrm{ind}}(x)/\mu(x)^2\le N^\alpha-1
\).
The bound is scale-invariant under common rescaling of the $X_i$.
\end{proposition}

\begin{proof}[Proof sketch]
Established in \cref{appx:independent-proofs} via the deterministic envelope
$Z^{\mathrm{ind}}(x)\le\sum_i\Fb_i(x/N)$ and the regular-variation limit
$\sum_i\Fb_i(x/N)/\mu(x)\to N^\alpha$.
\end{proof}

\begin{proposition}[BRE of dependent CdMC under a conditional-tail envelope]
\label{prop:dep-cdmc-bre}
Under \cref{ass:conditional-envelope}, the dependent CdMC estimator
$Z^{\mathrm{dep}}(x)$ is unbiased for $\mu(x)$ and satisfies
\(\limsup_{x\to\infty}\Var Z^{\mathrm{dep}}(x)/\mu(x)^2<\infty\).
\end{proposition}

\begin{proof}[Proof sketch]
This is \cref{prop:dep-envelope-bre}; see \cref{appx:dependent-proofs}.
\end{proof}

\begin{proposition}[BRE of latent-shock CdMC]
\label{prop:latent-shock-bre}
If $F=BZ$ with $Z_1,\ldots,Z_K$ independent and regularly varying with common
index, and \cref{ass:ind-sbj} holds for the signed shock contributions
$q_kZ_k$, then the latent-shock CdMC estimator is unbiased and satisfies the
independent BRE bound $K^\alpha-1$ in place of $N^\alpha-1$.
\end{proposition}

\begin{proof}[Proof sketch]
This is \cref{thm:latent-shock-tail} combined with the BRE bound for the
independent CdMC applied to the shock vector. The full argument is in
\cref{appx:dependent-proofs}.
\end{proof}

\begin{proposition}[BRE of spectral CdMC]
\label{prop:spectral-cdmc-bre}
Under \cref{ass:mrv-loss} with $\E[\ell(\Theta)_+^{2\alpha}]<\infty$ and the
envelope condition of \cref{prop:spectral-bre}, the spectral/radial CdMC
estimator $Z^{\mathrm{spec}}(x)$ is unbiased and has bounded relative error.
\end{proposition}

\begin{proof}[Proof sketch]
\Cref{prop:spectral-bre}; full argument in \cref{appx:dependent-proofs}.
\end{proof}

\subsection{Control variates and the VRE estimator}

The independent second-order approximation and the MRV spectral approximation
both produce analytic or semi-analytic controls. If \(Y(x)\) has known mean
\(m_Y(x)\), then
\[
  Z^{\mathrm{cv}}(x)=Z(x)-\gamma\{Y(x)-m_Y(x)\}
\]
is unbiased. The oracle coefficient is
\[
  \gamma^*(x)=\frac{\Cov(Z(x),Y(x))}{\Var(Y(x))}.
\]
In practice \(\gamma\) and sometimes \(m_Y\) are estimated on a pilot split.

\begin{proposition}[Exact and asymptotic VRE for the centered control-variate estimator]
\label{prop:vre}
Let \(Y(x)\) be an $L^2$ surrogate for \(Z(x)\) with $\Var Z(x),\Var Y(x)>0$
and let
\(\rho(x)=\mathrm{Corr}(Z(x),Y(x))\). Then:
\begin{enumerate}
\item \emph{(Exact VRE with oracle centering.)} If \(m_Y(x)\) is known exactly
and the oracle coefficient \(\gamma^*(x)\) is used, then $Z^{\mathrm{cv}}$ is
unbiased and
\[
  \Var Z^{\mathrm{cv}}(x)=(1-\rho(x)^2)\Var Z(x).
\]
If $\rho(x)\to 1$ as $x\to\infty$, then $\Var Z^{\mathrm{cv}}(x)/\mu(x)^2\to 0$
provided $Z(x)$ has BRE; that is, $Z^{\mathrm{cv}}$ has VRE.
\item \emph{(Asymptotic VRE with sample-split centering.)} If \(m_Y(x)\) is
estimated on a pilot split of size $n_0$ and \(\gamma\) on a second split, then
$Z^{\mathrm{cv}}$ is unbiased as long as the splits are independent of the
production sample, but its variance picks up an additional term of order
$1/n_0$ from the centering estimate. Under the same conditions as part (1)
and provided $n_0\to\infty$ at any rate with $x$, the work-normalized variance
ratio $W(Z^{\mathrm{cv}};x)/\Var Z(x)\to 0$ along the same threshold sequence.
\end{enumerate}
\end{proposition}

\begin{proof}[Proof sketch]
Part (1) is the standard identity for an oracle-centered control variate; the
VRE implication follows because $1-\rho(x)^2\to 0$ kills the surviving BRE
constant. Part (2) is the standard sample-splitting argument with an
independent pilot. Details mirror the classical control-variate analysis in
\citet{Glasserman2004}; see \cref{appx:dependent-proofs} for the bookkeeping.
\end{proof}

The distinction in \cref{prop:vre} matters for empirical claims: exact
centering yields VRE, while sample-split centering yields VRE only along
the joint limit $x\to\infty$, $n_0\to\infty$. The paper reports both variants.

\subsection{Finite-sample concentration}

If an estimator is bounded by \(0\le Z(x)\le B_Z(x)\), Hoeffding gives
\[
  \Pb\!\left(|\bar Z_n-\mu(x)|>\varepsilon\mu(x)\right)
  \le 2\exp\!\left\{-\frac{2n\varepsilon^2\mu(x)^2}{B_Z(x)^2}\right\}.
\]
A tighter Bernstein-style bound that exploits the small variance of CdMC
estimators is the following.

\begin{theorem}[Bernstein-type confidence interval for bounded CdMC estimators]
\label{thm:bernstein-ci}
Let $Z(x)$ be an unbiased estimator of $\mu(x)$ with $0\le Z(x)\le B_Z(x)$ and
$\Var Z(x)\le \sigma^2(x)$. Let $\bar Z_n$ denote the sample mean of $n$ i.i.d.\
replicates. Then for all $\varepsilon>0$,
\[
  \Pb\!\left(|\bar Z_n-\mu(x)|>\varepsilon\right)
  \le 2\exp\!\left\{-\frac{n\varepsilon^2}{2\sigma^2(x)+\tfrac{2}{3}B_Z(x)\varepsilon}\right\}.
\]
In particular, taking $\varepsilon=\delta\mu(x)$ and using the BRE bound
$\sigma^2(x)/\mu(x)^2\le V$ and the envelope ratio
$B_Z(x)/\mu(x)\le M$ gives
\[
  \Pb\!\left(|\bar Z_n-\mu(x)|>\delta\mu(x)\right)
  \le 2\exp\!\left\{-\frac{n\delta^2}{2V+\tfrac{2}{3}M\delta}\right\}.
\]
\end{theorem}

\begin{proof}[Proof sketch]
This is Bernstein's inequality applied to the i.i.d.\ sequence
$Z_1(x),\ldots,Z_n(x)$, with the variance and uniform bound substituted in.
The substitution to relative-error form uses the BRE bound on $\sigma^2/\mu^2$
and the deterministic envelope ratio $B_Z/\mu$; for the independent CdMC,
$M=N^\alpha$ and $V=N^\alpha-1$ by \cref{prop:ind-cdmc-bre}.
\end{proof}

The simulation study reports Bernstein intervals from \cref{thm:bernstein-ci},
Hoeffding intervals when only a deterministic range bound is available, and
self-normalized confidence sequences when heavy-tailed estimator outputs make
a range bound too conservative.

\subsection{Planned estimator workflow}

The workflow is shown in \cref{fig:estimator-workflow}. It deliberately keeps
estimation and validation separate: tail and dependence diagnostics determine
candidate estimators; out-of-sample VaR/ES backtests determine whether the
added dependence complexity improves risk forecasts.

\begin{figure}[p]
\centering
\begin{tikzpicture}[
  box/.style={rectangle, draw, rounded corners, align=center, minimum width=2.55cm, minimum height=0.8cm, font=\scriptsize},
  node distance=0.55cm and 0.55cm,
  arr/.style={->, >=Stealth, thick}
]
\node[box, fill=blue!10] (loss) {Loss contribution\\matrix};
\node[box, fill=blue!10, right=of loss] (marg) {Marginal tail\\fit};
\node[box, fill=blue!10, right=of marg] (dep) {Dependence\\diagnostics};
\node[box, fill=yellow!15, below=of dep] (select) {Estimator\\selection};
\node[box, fill=green!12, left=of select] (rare) {Rare-event\\estimation};
\node[box, fill=green!12, left=of rare] (risk) {VaR/ES\\forecast};
\node[box, fill=orange!12, below=of risk] (back) {Backtests and\\attribution};
\draw[arr] (loss) -- (marg);
\draw[arr] (marg) -- (dep);
\draw[arr] (dep) -- (select);
\draw[arr] (select) -- (rare);
\draw[arr] (rare) -- (risk);
\draw[arr] (risk) -- (back);
\draw[arr] (back.east) -- ++(1.0,0) |- (dep.south);
\end{tikzpicture}

\caption[Estimator workflow]{Estimator workflow from raw loss contributions to
model diagnostics, rare-event estimation, VaR/ES forecasts, and backtests. The
same workflow is used for synthetic experiments and for real data.}
\label{fig:estimator-workflow}
\end{figure}
